% Title page, table of contents and "About This Book".

% The title page and the first TOC page are both page 1; without this, hyperref
% creates the anchor page.1 twice (as in algebra-1.tex).
\hypersetup{pageanchor=false}
\begin{titlepage}
\thispagestyle{empty}
\begin{tikzpicture}[remember picture, overlay]
  \fill[ThemeSecondary!10, rounded corners=8pt] ($(current page.north east) + (-7.2cm,-2.2cm)$) rectangle ($(current page.north east) + (-1.2cm,-6.8cm)$);
  \fill[ThemePrimary!12, rounded corners=8pt] ($(current page.north east) + (-5.2cm,-3.6cm)$) rectangle ($(current page.north east) + (-0.6cm,-8.2cm)$);
  \fill[TextBoldColor!15, rounded corners=6pt] ($(current page.north east) + (-3.4cm,-5.0cm)$) rectangle ($(current page.north east) + (-0.6cm,-9.4cm)$);
\end{tikzpicture}

\vspace*{4.2cm}
\noindent\begin{minipage}{\textwidth}
    {\sffamily\fontsize{42pt}{50pt}\selectfont\bfseries\color{ThemePrimary!95!black} Algebra I \par}
    \vspace{0.8cm}
    {\noindent\textcolor{ThemePrimary}{\rule{\textwidth}{1.5pt}}\par}%
    \vspace{1.8cm}
    {\sffamily\fontsize{18pt}{24pt}\selectfont\bfseries\color{ThemeSecondary!90!black} A Book Based on the Lectures \par}
    \vspace{0.8cm}
    {\sffamily\fontsize{13pt}{22pt}\selectfont
      \begin{tabular}{@{}ll@{}}
        \textcolor{black!50}{Lectures:}    & \textbf{\textcolor{ThemeSecondary!90!black}{Prof.\ Rahul Pandharipande}} \\[0.5ex]
        \textcolor{black!50}{Institution:} & \textcolor{black!75}{Department of Mathematics -- ETH Z{\"u}rich} \\[0.5ex]
        \textcolor{black!50}{Semester:}    & \textcolor{ThemeWeekNumber}{\faGraduationCap\quad Fall Semester 2026} \\[0.5ex]
        \textcolor{black!50}{Edition:}     & \textcolor{black!75}{Pilot: Chapter 6 only}
      \end{tabular}\par}
\end{minipage}
\vfill
\noindent\begin{minipage}{\textwidth}
    \begin{tcolorbox}[
        arc=6pt,
        boxrule=0.8pt,
        colback=MidnightBlue!2!white,
        colframe=MidnightBlue!25!white,
        left=12pt, right=12pt, top=10pt, bottom=10pt
    ]
    {\sffamily\bfseries\small\color{MidnightBlue!90!black}\faInfoCircle\quad Notice \par}
    \vspace{0.25cm}
    \sffamily\small\color{black!75}
    \begin{itemize}[leftmargin=1.5em, itemsep=3pt, topsep=0pt, parsep=0pt]
        \item[\textcolor{MidnightBlue!80!black}{\small\faRobot}] \textbf{Written with AI:} This book was written with AI systems from the lecture notes \emph{Algebra I -- Lecture Notes \& Transcription}, which were themselves transcribed by AI from the recordings and checked against them.
        \item[\textcolor{MidnightBlue!80!black}{\small\faUserShield}] \textbf{Unofficial Document:} This is \emph{not} official course material of ETH Z{\"u}rich or of the lecturer, and it has \emph{not} been read or corrected by the course staff.
        \item[\textcolor{MidnightBlue!80!black}{\small\faExclamationTriangle}] \textbf{Disclaimer:} AI can make errors (in notation, indices or logical steps). Use at your own risk, without warranty of any kind regarding accuracy, completeness, or relevance to exams.
    \end{itemize}
    \end{tcolorbox}
\end{minipage}
\vspace*{0.5cm}
\end{titlepage}
\hypersetup{pageanchor=true}

\tableofcontents
\hypersetup{linkcolor=LinkColor}
\cleardoublepage

\chapter*{About This Book}
\addcontentsline{toc}{chapter}{About This Book}
\markboth{About This Book}{}

This book presents the course Algebra~I as Prof.\ Rahul Pandharipande taught it at ETH Z{\"u}rich in the fall semester of 2026, ordered by topic rather than by date. Its source is the lecture notes \emph{Algebra I -- Lecture Notes \& Transcription}, a record of every lecture with what was said and what was written on the board, checked against the recordings. Those notes stay the place to look up a lecture as it happened.

The book keeps the mathematics of the lectures: the definitions, statements and proofs from the board, and the lecturer's notation. The spoken explanations were used as a source for the motivation and the pictures that come with the mathematics, but they are not reproduced. A topic that the lectures treated more than once, for example a definition restated at the start of the next lecture, appears once.

Two kinds of additions are marked:
\begin{itemize}
    \item mathematics the lectures did not give, such as a proof stated as an exercise or a standard fact, stands in a box headed \emph{Beyond the Lecture};
    \item names and notations that textbooks use but the lectures did not are given in footnotes.
\end{itemize}

\section*{The Pilot Edition}

Only Chapter~6 is written so far. Chapters 0 to 5 are placeholders: each holds just the statements that Chapter~6 refers to, copied from the lecture notes, so that its cross-references point to the right place.

\section*{Conventions}

\begin{itemize}
    \item \textbf{Rings} are commutative with a unit $1$, and the zero ring is allowed (\cref{rem:ring-convention}).
    \item Definitions are written with $:=$ or with $\stackrel{\text{def}}{\iff}$, as on the board.
    \item $A \subset B$ allows $A = B$; a strict inclusion is written $A \subsetneq B$ (\cref{not:subset-board}).
    \item Ideals are written with plain letters: $I$, $J$, $K$, and $m$ for a maximal ideal. Sets of ideals are written with a script $\mathscr{I}$ (\cref{not:ideal-sets-k}, \cref{def:proper-ideals}).
\end{itemize}
